\documentclass[10pt,a4paper]{amsart}
\usepackage[margin=19mm]{geometry}
\usepackage{amsmath,amssymb,mathtools}
\usepackage[scr=rsfs,cal=euler]{mathalfa}
\usepackage{xfrac}
\usepackage[unicode,hidelinks,bookmarks=false]{hyperref}
\newcommand{\ie}{i.e.\ }
\newcommand{\cf}{cf.\ }
\newcommand{\resp}{respectively\ }
\newcommand{\Subsection}{\subsection}
\providecommand{\nobreakdash}{\nobreak\hbox{-}}
\setlength{\emergencystretch}{2em}
\multlinegap0pt
\usepackage{mathtools}
\usepackage{amssymb}
\usepackage[shortlabels]{enumitem}
\usepackage{cancel}
\usepackage{bm}
\usepackage{xcolor}
\newtheorem{lemma}{Lemma}
\title{Scattering for the quintic generalized Benjamin-Bona-Mahony Equation}
\author{Gong Chen, Yingmo Zhang}
\date{Source: arXiv:2603.01360v1 (CC BY 4.0)}
\begin{document}
\maketitle

\noindent\emph{Source and licence.} Scattering for the quintic generalized Benjamin-Bona-Mahony Equation. Version arXiv:2603.01360v1 (CC BY 4.0), distributed by arXiv under the Creative Commons Attribution 4.0 International licence, http://creativecommons.org/licenses/by/4.0/, as recorded in the arXiv licence metadata for this exact version and shown on the abstract page https://arxiv.org/abs/2603.01360v1. Full licence notice: \texttt{licenses/arxiv-2603.01360-CC-BY-4.0.txt}.\par\smallskip

\noindent\textit{Editorial scope.} Counted unit: the article's lemma lem:anomalous\_res (source0.tex lines 504--506). The article's complete original proof of it is delivered with it (source0.tex lines 507--606, one \texttt{proof} environment of 100 lines). No mathematics is written, repaired or re-derived: the statement and the proof are the article's own lines, in the article's own notation, and no retained line is substituted or re-flowed.  Retained uncounted as the source-local context the proof needs: lines 404--406.  Nothing was omitted from any retained span, so the package carries no omission record.  No retained line contains a citation, so the wrapper carries no bibliography and no bibliography entry.  No retained line contains a figure environment, an image-inclusion command or drawing code, so no image is delivered and the package declares no asset dependency.  Editorial formatting only, in addition: an AMS \texttt{amsart} preamble that restates the article's own declarations and macros used by the retained lines, namely the environments \texttt{lemma} on separate counters, together with the standard packages those lines need.  The retained environments print in this document as Lemma 1 in order of appearance, whereas the article prints the same items with the numbering of its own full document.  The complete native source of the article as distributed in the arXiv e-print for this exact version is bundled unchanged as \texttt{sources/195555/source0.tex}.
    \begin{equation}\label{eq:def_xi_eta_nod}
        \eta_0\approx 7.34 \qquad \text{and}\qquad \xi_0\approx 28.31
    \end{equation}

\begin{lemma}\label{lem:anomalous_res}
    For the coefficient pairs $(A,B) \in \{(4,-1), (-4,1)\}$, the equation $\Phi_{A,B}(\eta) = 0$ admits a unique positive root $\eta_0 \approx 7.34$ and its negative counterpart $-\eta_0$. This root corresponds to the anomalous resonance defined in \eqref{eq:def_xi_eta_nod}.
\end{lemma}

\begin{proof}
    By the odd symmetry of $\omega$, it suffices to consider $(A,B) = (4,-1)$. The equation becomes $4\omega(\eta) - \omega(r(\eta)) - \omega(4\eta - r(\eta)) = 0$.

Denote 
\begin{equation}
    H(\eta):=4\omega(\eta)-\omega\big(r(\eta)\big)
-\omega\big(4\eta-r(\eta)\big).
\end{equation}Note that $H$ is an odd function so it suffices to consider its zeros for  $1<\eta$.

We record some elementary computations. For $\eta>1$ we have
$$
r(\eta)=\sqrt{\frac{\eta^2+3}{\eta^2-1}}>1.
$$
A direct computation gives
$$
1+r(\eta)^2=\frac{2(\eta^2+1)}{\eta^2-1},
\qquad
\omega\big(r(\eta)\big)
=\frac{\sqrt{(\eta^2+3)(\eta^2-1)}}{2(\eta^2+1)}.
$$
Moreover
$$
\omega'(\xi)=\frac{1-\xi^2}{(1+\xi^2)^2},
\qquad
|\omega(\xi)|\le \frac12 \quad\text{for all }\xi.
$$
Differentiating $r$ for $\eta>1$,
%$$
%\frac{r'(\eta)}{r(\eta)}
%=\frac12\Big(\frac{2\eta}{\eta^2+3}-%\frac{2\eta}{\eta^2-1}\Big)
%=-\frac{4\eta}{(\eta^2+3)(\eta^2-1)},
%$$
%hence
\begin{equation}\label{eq:rprime}
r'(\eta)=-\frac{4\eta\,r(\eta)}{(\eta^2+3)(\eta^2-1)}<0.
\end{equation}
For $1<\eta\le2$, with $|\omega|\le\frac12$, one has
$$
\omega\big(r(\eta)\big)+\omega\big(4\eta-r(\eta)\big)\le1.
$$
On the other hand,
$$
4\omega(\eta)=\frac{4\eta}{1+\eta^2}
\ge \frac{8}{5}>1,
$$
where the minimum on $[1,2]$ occurs at $\eta=2$. Thus
$$
H(\eta)>0\qquad\text{for }1<\eta\le2.
$$
For $\eta\ge2$, one has
$$
H'(\eta)
=4\omega'(\eta)
-\omega'\big(r(\eta)\big)\,r'(\eta)
-\omega'\big(4\eta-r(\eta)\big)\,(4-r'(\eta)).
$$
Since $r(\eta)>1$ and $r'(\eta)<0$, we have $
-\omega'(r(\eta))\,r'(\eta)<0$,
so it suffices to bound the remaining terms from above.    

For $\eta\ge2$,
$$
4\omega'(\eta)
=\frac{4(1-\eta^2)}{(1+\eta^2)^2}
\le -\frac{1}{2\eta^2}.
$$
Also $x(\eta):=4\eta-r(\eta)>1$ and
$$
-\omega'(x)\le \frac{1}{x^2}.
$$
From \eqref{eq:rprime}, $|r'(\eta)|\le1$ for $\eta\ge2$, hence
$4-r'(\eta)\le5$. Therefore
$$
-\omega'(x(\eta))\,(4-r'(\eta))
\le \frac{5}{(4\eta-r(\eta))^2}
\le \frac{5}{(4\eta-\sqrt{\frac{7}{3}})^2}.
$$
A direct comparison shows that for all $\eta\ge2$,
$$
\frac{5}{(4\eta-\sqrt{\frac{7}{3}})^2}
<\frac{1}{2\eta^2},
$$
and consequently $H'(\eta)<0$ for $\eta\ge2$. Thus $H$ is strictly
decreasing on $[2,\infty)$.

We already have $H(2)>0$. Moreover, as $\eta\to\infty$,
$$
r(\eta)\to1,\qquad
\omega\big(r(\eta)\big)\to\frac12,\qquad
\omega(4\eta-r(\eta))\to0,\qquad
4\omega(\eta)\to0,
$$
so
$$
\lim_{\eta\to\infty}H(\eta)=-\frac12<0.
$$
By continuity and strict monotonicity on $[2,\infty)$, there exists a
unique $\eta_0>2$ such that $H(\eta_0)=0$.
Moreover \eqref{eq:def_xi_eta_nod} is the only solution\footnote{Verified by Mathematica}.
\end{proof}

\end{document}
